% Generated draft: requires mathematical and source review.
\documentclass[11pt]{article}
\usepackage[margin=25mm]{geometry}
\usepackage{fontspec}
\setmainfont{Latin Modern Roman}
\usepackage{amsmath,amssymb,mathtools,unicode-math}
\setmathfont{Latin Modern Math}
\usepackage{enumitem,needspace,xurl,hyperref}
\setlength{\parindent}{0pt}
\setlength{\parskip}{5pt}
\setlength{\emergencystretch}{3em}
\newcommand{\N}{\mathbb{N}}
\newcommand{\Z}{\mathbb{Z}}
\newcommand{\Q}{\mathbb{Q}}
\newcommand{\R}{\mathbb{R}}
\newcommand{\C}{\mathbb{C}}
\newcommand{\F}{\mathbb{F}}
\newcommand{\A}{\mathbb{A}}
\newcommand{\PP}{\mathbb{P}}
\newcommand{\E}{\mathbb{E}}
\newcommand{\eps}{\varepsilon}
\newcommand{\dd}{\,\mathrm{d}}
\begin{document}
\section*{UnsolvedMath ID: 700577}
\textbf{Source problem code:} catalogue entry
\subsection{Definitions and mathematical statement}
Let \(G\) be a topological group. For every integer \(n\geq 1\), define the power map \(p_n:G\to G\) by \(p_n(x)=x^n\). Following the source, call \(G\) \emph{layer compact} if, for every compact subset \(K\subseteq G\) and every integer \(n\geq 1\), the full inverse image
\[
 p_n^{-1}(K)=\{x\in G:x^n\in K\}
\]
is compact in \(G\).

Use the usual group-theoretic meaning of \emph{locally soluble}: every finitely generated subgroup \(H\leq G\) is soluble. Explicitly, with \(H^{(0)}=H\) and \(H^{(j+1)}=[H^{(j)},H^{(j)}]\), for each such \(H\) there exists an integer \(d\geq 0\) such that \(H^{(d)}=\{1\}\). The integer \(d\) may depend on \(H\); no uniform bound is imposed, and these are abstract subgroups, without taking topological closures.

\textbf{Problem.} Describe all locally compact, locally soluble topological groups \(G\) that are layer compact. Here local compactness means that every point has a compact neighbourhood. The supplied record does not specify a separation convention for topological groups, in particular whether Hausdorffness is assumed; this convention requires clarification. The interpretation of local solubility above is the standard algebraic reading, but the supplied record gives no explicit definition.
\subsection{Short English statement}
Describe the locally compact groups in which every finitely generated subgroup is soluble and, for every positive integer \(n\), the set of all \(n\)-th roots of elements of any compact subset is compact. The source leaves the separation convention unspecified. Its label \emph{open} is supplied registry information, not an independently verified assessment.
\subsection{Sources}
\begin{itemize}
\item Local UnsolvedMath registry, permanent numeric ID \(700577\), exact source problem code \texttt{catalogue entry}, title \emph{Kourovka Notebook Problem 9.44}. The supplied record labels the problem \emph{open}.
\item \emph{The Kourovka Notebook (Unsolved Problems in Group Theory)}, Problem 9.44, proposed by V. M. Poletskikh; 9th issue (1984). The supplied background also identifies the 21st issue edition (2026), \texttt{arXiv:1401.0300v48}. The supplied collection description names E. I. Khukhro and V. D. Mazurov as editors of Issues 1--21 (1965--2026). These bibliographic details are supplied attribution and have not been independently checked.
\end{itemize}
\end{document}
